\section{Entrevista a un Operador de Cafetalino}\label{anexo-entrevista-a-un-operador-de-cafetalino}

\textbf{GUÍA PARA LA ENTREVISTA A UN OPERADOR DE RECARGA DE LA RED DE MÁQUINAS DISPENSADORAS DE CAFETALINO}

El objetivo de la entrevista es conocer y comprender el proceso de recarga de insumos a la red de máquinas dispensadoras de Cafetalino, los problemas y particularidades del mismo.

\begin{enumerate}
\def\labelenumi{\arabic{enumi}.}
\item
  ¿Cuál es el procedimiento de recarga de las máquinas dispensadoras?
\item
  ¿Cómo se definen las rutas de recarga?
\item
  ¿Alguna vez ha tenido que modificar la ruta de recarga? ¿Por qué?
\item
  ¿Cuán eficiente considera que es este procedimiento?
\item
  ¿Consideraría útil un sistema que le ayude a priorizar sus rutas en función a la demanda de recarga de las máquinas de la red de Cafetalino?
\end{enumerate}

\textbf{DESARROLLO DE LA ENTREVISTA A UN OPERADOR DE RECARGA DE LA RED DE MÁQUINAS DISPENSADORAS DE CAFETALINO}

\begin{enumerate}
\def\labelenumi{\arabic{enumi}.}
\item
  ¿Cuál es el procedimiento de recarga de las máquinas dispensadoras?
\end{enumerate}

\begin{quote}
\emph{Cada día recorremos dos rutas que comprenden entre 4 y 6 maquinas, una en la mañana y otra en la tarde.}

\emph{Antes de salir preparamos los paquetes de recarga de insumos, uno por maquina y 2 extra por si es necesario. Luego hacemos el recorrido y en cada ubicación, apagamos la máquina, la revisamos, le hacemos limpieza y un mantenimiento rutinario, la recargamos y la encendemos nuevamente para ponerla en servicio, verificando su correcto funcionamiento, que no haya fugas, ruidos extraños, ni alertas de ningún tipo.}
\end{quote}

\begin{enumerate}
\def\labelenumi{\arabic{enumi}.}
\setcounter{enumi}{1}
\item
  ¿Cómo se definen las rutas de recarga?
\end{enumerate}

\begin{quote}
\emph{Las rutas están predefinidas, cada ruta cubre un grupo de máquinas que están más o menos por la misma zona, el número de máquinas está en función a la distancia entre maquinas.}
\end{quote}

\begin{enumerate}
\def\labelenumi{\arabic{enumi}.}
\setcounter{enumi}{2}
\item
  ¿Alguna vez ha tenido que modificar la ruta de recarga? ¿Por qué?
\end{enumerate}

\begin{quote}
\emph{Si, a veces, nos avisan que hay una máquina fuera de servicio, si está en la misma ruta, le damos prioridad y vamos primero a esa ubicación, y luego cambiamos la secuencia para no estar yendo de un lado a otro, lo grave es si es de otra ruta, porque siempre hay que priorizar la que se quedó fuera de servicio y luego volver a la ruta normal o invertir la secuencia para no estar yendo de un lado a otro, para estos casos es que preparamos 2 paquetes de recarga extra.}
\end{quote}

\begin{enumerate}
\def\labelenumi{\arabic{enumi}.}
\setcounter{enumi}{3}
\item
  ¿Cuán eficiente considera que es este procedimiento?
\end{enumerate}

\begin{quote}
\emph{No es muy eficiente, porque como le dije a veces es necesario cambiar la secuencia de la ruta o hasta ir a otras rutas y muchas veces hay que hacer horas extra para terminar de revisar todas las maquinas. Además, a veces llegamos a una ubicación que por algún motivo no ha tenido el consumo regular, porque hubo corte de energía eléctrica en la zona o alguna otra cosa, lo que pasa en estos casos es que no requiere recarga y hacerle el mantenimiento se complica porque los depósitos están llenos, y bueno solo se la limpia y se la pone en servicio nuevamente, pero es solo una pérdida de tiempo}
\end{quote}

\begin{enumerate}
\def\labelenumi{\arabic{enumi}.}
\setcounter{enumi}{4}
\item
  ¿Consideraría útil un sistema que le ayude a priorizar sus rutas en función a la predicción de la demanda de recarga de las máquinas de la red de Cafetalino?
\end{enumerate}

\begin{quote}
\emph{Uhh, sería una gran cosa, pues se evitarían estas pérdidas de tiempo, la ruta se programaría con las maquinas que realmente requieren recarga, y no se tendría que modificar las rutas a medio recorrido, o quizás sí, pero sería excepcionalmente.}
\end{quote}
